\scanpage{99}
\mspage{86}

\deleted{La première des relations (5.6) est une tautologie, la deuxième \deleted{pour} équivaut au fait que $\mathrm{TC}$ et $\mathrm T\widetilde C$ ont $\mathord{=}$ Q-clôture à gauche, déjà prouvé (5.5.).}

Cela achève de prouver b). Pour b'), il reste à voir que $\mathrm T\widetilde C$ est stable par composition transfinie. Je ne vois pas d'autre façon de le voir, qu'en supposant que $W$ lui-$=$ est stable par composition transfinie (ce qui paraît à peine moins exigeant que la stabilité de $W$ par $\varinjlim$). \inserted{Cela achève de prouver b,b').} La condition c) sur $C$ \deleted{\unsure{signi|sign}} implique la $=$ condition sur $\widetilde C$. Enfin, comme $\widetilde C_{\unsure{*|\alpha}}=C_*$, la condition d) est la $=$ pour $C$ et $\widetilde C$. Cela achève de prouver que $(\unsure{\mathcal M|M},W,\widetilde C)$ est une cat. à cofibrations spéciale. (\uline{NB} On avait besoin des \deleted{\unsure{2 hyp|2 hy}} deux hyp. suppl. sur $W$, stabilité par facteur direct et par composition transfinie, pour établir que $\mathrm T\widetilde C$ est stable par cochg\textsuperscript{t} de base et par composition transfinie). La dernière assertion de 5.1 \unsure{de|et} la relation (5.7) qui en résulte, \unsure{revient} à dire que $\mathrm T\widetilde C=\widetilde C\cap W$ est stable par facteurs directs, ce qui est évident par la $=$ \deleted{\unsure{hyp.}} propriété de $\widetilde C$, et l'hyp. faite sur $W$ de \deleted{que les facteurs directs d'\unsure{élém}} stabilité par fact. directs, (ou seulement que les fact. directs d'éléments de $\mathrm T\widetilde C$ sont dans $W$, donc dans $\mathrm T\widetilde C$).
% archive sheet number: 98
